working at the categorical rather than algebraic
level. Concretely, the Drinfeld centre
$\cZ(Y_\hbar(\mathfrak{sl}_2)^{\mathrm{ch}}\textrm{-mod})$
should recover the quantum group category
$\mathrm{Rep}_q(\mathfrak{sl}_2)$ at
$q = e^{2\pi i/(k+2)}$. This identification is the
subject of the DK-5 programme; see
\cite[\S6]{Lorgat26I} for the current status in type~$A$.
\end{remark}

The computations in
Example~\ref{ex:yangian} and
Conjecture~\ref{conj:yangian-ordered-center}(v)
concern the kernel of the averaging map on the
ordered chiral centre, a cohomological object.
The following proposition establishes the underlying
representation-theoretic input: the dimension of
$\ker(\av_n)$ on the bare tensor power $V^{\otimes n}$,
before passing to bar-complex cohomology.

\begin{proposition}[Kernel of the Reynolds projector:
general simple Lie algebras]
\label{prop:ker-av-schur-weyl}
Let $\fg$ be any simple Lie algebra over~$\CC$, let
$V$ be any finite-dimensional representation of~$\fg$ with
$d = \dim V$, and let
$\av_n \colon V^{\otimes n} \to (V^{\otimes n})_{\Sigma_n}$
denote the Reynolds averaging map (projection onto the
$\Sigma_n$-coinvariant quotient).
Then
\begin{equation}\label{eq:ker-av-schur-weyl}
  \dim \ker(\av_n\big|_{V^{\otimes n}})
  \;=\;
  d^n - \binom{n+d-1}{d-1}.
\end{equation}
The formula depends only on $d = \dim V$, not on $\fg$ or
on the $\fg$-module structure of~$V$.

\smallskip
\noindent\textit{Special cases.}
For $\fg = \mathfrak{sl}_2$ with the fundamental
representation $V = \CC^2$:
$\dim\ker(\av_n) = 2^n - (n+1)$.
For any $\fg$ with $d$-dimensional representation~$V$,
at $n = 2$:
$\dim\ker(\av_2) = d^2 - d(d{+}1)/2 = d(d{-}1)/2
= \dim\bigwedge^2 V$.
\end{proposition}

\begin{proof}
The Reynolds projector
$\av_n = \frac{1}{n!}\sum_{\sigma \in \Sigma_n} \sigma$
acts on $V^{\otimes n}$ by permuting tensor factors.
Its image is the $\Sigma_n$-invariant subspace
$\Sym^n(V)$, regardless of any Lie algebra structure
on~$V$. The dimension
$\dim\Sym^n(V) = \binom{n+d-1}{d-1}$ (the number of
degree-$n$ monomials in $d$ variables) depends only on
$d = \dim V$, as the stars-and-bars argument uses no
algebraic structure beyond the vector space.
Since $\dim V^{\otimes n} = d^n$, the kernel has dimension
$d^n - \binom{n+d-1}{d-1}$.

When $\fg = \mathfrak{gl}_r$ with $V = \CC^r$,
Schur--Weyl duality refines this to an isotypic
decomposition. The tensor power decomposes into
$(\Sigma_n, \mathfrak{gl}_r)$-bimodules:
\begin{equation}\label{eq:schur-weyl-decomp}
  V^{\otimes n}
  \;\cong\;
  \bigoplus_{\lambda \,\vdash\, n,\;
    \ell(\lambda) \leq r}
  S_\lambda \otimes V_\lambda(\mathfrak{gl}_r),
\end{equation}
where $\lambda$ runs over partitions of $n$ with at most
$r$ rows, $S_\lambda$ is the corresponding Specht module
of $\Sigma_n$, and $V_\lambda(\mathfrak{gl}_r)$ is the
irreducible $\mathfrak{gl}_r$-representation of highest
weight~$\lambda$.
The projector $\av_n$ selects the trivial isotypic
component $S_{(n)} = \CC$ and annihilates every
$S_\lambda$ with $\lambda \neq (n)$. The image is
\[
  \operatorname{im}(\av_n)
  \;\cong\;
  S_{(n)} \otimes V_{(n)}(\mathfrak{gl}_r)
  \;=\;
  \Sym^n(\CC^r),
\]
and the kernel decomposes by isotypic components:
\[
  \ker(\av_n)
  \;\cong\;
  \bigoplus_{\substack{\lambda \,\vdash\, n \\
    \ell(\lambda) \leq r,\;
    \lambda \neq (n)}}
  S_\lambda \otimes V_\lambda(\mathfrak{gl}_r).
\]
For a general simple~$\fg$ acting on~$V$, the Schur--Weyl
decomposition involves the commutant algebra
$\End_\fg(V^{\otimes n})$ rather than the group algebra of
$\GL(V)$, and the individual summands carry different
$\fg$-representations with different dimensions.
The total kernel dimension, however, depends only on
$\dim\Sym^n(V) = \binom{n+d-1}{d-1}$, which is
determined by $d$ alone.
\end{proof}

\begin{remark}[Explicit values and
Conjecture~\textup{\ref{conj:yangian-ordered-center}}]
\label{rem:ker-av-table}
Table~\ref{tab:ker-av-dims} records the kernel dimensions
at low degrees for the fundamental representations of
$\mathfrak{sl}_2$ through $\mathfrak{sl}_5$ (of dimensions
$d = 2, 3, 4, 5$), together with the fundamental
representations of $\mathfrak{so}_5$ ($d = 5$) and
$G_2$ ($d = 7$). Since the formula depends only on
$d = \dim V$, the rows for $\mathfrak{sl}_5$ and
$\mathfrak{so}_5$ are identical.
These are representation-theoretic upper bounds on the
dimensions of the corresponding components of the ordered
chiral centre in
Conjecture~\ref{conj:yangian-ordered-center}(v):
$\ker(\av_n)$ on the ordered chiral centre is a subquotient
of $\ker(\av_n\big|_{V^{\otimes n}})$, obtained by passing
to bar-complex cohomology
(only the cocycles in the non-trivial $\Sigma_n$-isotypics
survive). At degree $2$ for $\mathfrak{sl}_2$: the full
$1$-dimensional kernel $\bigwedge^2 V = \CC$ survives
(Example~\ref{ex:yangian},
equation~\eqref{eq:yangian-ker-av2}).
At degree $n \geq 3$, the kernel on the chiral centre is
in general strictly smaller than the
representation-theoretic bound $d^n - \binom{n+d-1}{d-1}$.

\begin{table}[h]
\centering
\renewcommand{\arraystretch}{1.3}
\begin{tabular}{@{} l c c c c c @{}}
\toprule
  & $n = 2$ & $n = 3$ & $n = 4$ & $n = 5$ & $n = 6$ \\
\midrule
$\mathfrak{sl}_2$ $(d{=}2)$
  & $1$ & $4$ & $11$ & $26$ & $57$ \\
$\mathfrak{sl}_3$ $(d{=}3)$
  & $3$ & $17$ & $66$ & $222$ & $701$ \\
$\mathfrak{sl}_4$ $(d{=}4)$
  & $6$ & $44$ & $221$ & $968$ & $4{,}012$ \\
$\mathfrak{sl}_5$, $\mathfrak{so}_5$ $(d{=}5)$
  & $10$ & $90$ & $555$ & $2{,}999$ & $15{,}415$ \\
$G_2$ \textit{fund.}\ $(d{=}7)$
  & $21$ & $259$ & $2{,}191$ & $16{,}345$ & $116{,}725$ \\
\midrule
\textit{Formula}
  & \multicolumn{5}{c}{$d^n - \binom{n+d-1}{d-1}$} \\
\bottomrule
\end{tabular}
\caption{Dimension of $\ker(\av_n\big|_{V^{\otimes n}})$
for a $d$-dimensional representation $V$ at degrees
$n = 2, \ldots, 6$
\textup{(Proposition~\ref{prop:ker-av-schur-weyl})}.
The formula depends only on $d = \dim V$.}
\label{tab:ker-av-dims}
\end{table}
\end{remark}


% ================================================================
\section{Background: chiral algebras and the $\Eone/\Einf$
hierarchy}
\label{sec:background}
